\input{glyphtounicode}
\pdfgentounicode=1
\documentclass[11pt,a4paper]{article}
\usepackage{cmap}
\usepackage[utf8]{inputenc}
\usepackage[T1]{fontenc}
\usepackage[ngerman]{babel}
\usepackage{lmodern}
\usepackage{microtype}
\usepackage[a4paper,top=2.5cm,bottom=2.5cm,left=2.5cm,right=2.5cm]{geometry}
\usepackage{amsmath,amssymb,amsthm}
\usepackage{booktabs}
\usepackage{tabularx}
\usepackage[colorlinks=true,linkcolor=blue,citecolor=blue,urlcolor=blue,unicode=true]{hyperref}
\usepackage{xurl}

\renewcommand*{\HyperDestNameFilter}[1]{ger.#1}

\hypersetup{
  pdftitle={Die Joshi-ATS-Serie unter dem Brückenzertifikat-Standard: Ein gezielter Substanz-Audit und eine Allstellen-Starrheitsklassifikation der effektiven Normierungsgewichte},
  pdfauthor={Lukas Geiger},
  pdfsubject={Mathematik - Zahlentheorie, Arithmetische Geometrie},
  pdfkeywords={Arithmetischer Teichmüllerraum, abc-Vermutung, Brückenzertifikat, Substanz-Audit, Starrheitsklassifikation, Normierungsgewichte},
  pdfcreator={LaTeX with hyperref},
  pdfproducer={pdfTeX},
  bookmarksnumbered=true,
  pdfpagemode=UseOutlines,
  pdfstartview={FitH},
  unicode=true
}

\newtheorem{theorem}{Satz}[section]
\newtheorem{lemma}[theorem]{Lemma}
\newtheorem{proposition}[theorem]{Proposition}
\newtheorem{corollary}[theorem]{Korollar}
\theoremstyle{definition}
\newtheorem{definition}[theorem]{Definition}
\theoremstyle{remark}
\newtheorem{remark}[theorem]{Bemerkung}

\title{Die Joshi-ATS-Serie unter dem Brückenzertifikat-Standard:\\Ein gezielter Substanz-Audit und eine Allstellen-\\Starrheitsklassifikation der effektiven Normierungsgewichte}
\author{Lukas Geiger\\\small Unabhängiger Forscher, Bernau, Deutschland\\\small ORCID: 0009-0005-7296-1534}
\date{Oktober 2026 --- Version 0.3 (Entwurf)}

\begin{document}
\maketitle

\begin{abstract}
Dieser fortlaufende gezielte Audit trennt Koordinatenbuchführung,
Normierungsgewichte und Volumenträger an ausgewählten tragenden Schritten
der Joshi-ATS-Serie. Das Periodenargument verlangt konsistente rohe und
normierte Log-Koordinaten sowie eine wohldefinierte Aktion auf dem
Periodenbild. Gleichfeldrige Allstellen-Produktformel-Starrheit erzwingt
über die Stellen konstante Effektiv\-gewichte, erlaubt aber eine globale
$j^2$-Skalierung; ihre Anwendung über $L'/L$ und auf das Einfrieren der
Quelle bleibt bedingt. Eine lokale Maximalordnungs-Hülleninklusion gilt
unter einem tatsächlichen Modulhüllenvertrag; ATS-Logarithmus-, Haargrad-,
Tensor- und Zielvolumenbrücken bleiben offen. Das gedruckte ATS-III-v4-Korollar
besitzt Betragsstriche: Bei endlicher nichtleerer Tate-Stellensumme mit
$0<|q_w|<1$ und positiven $\ell,\ell^*$ ist links der Ausdruck
nichtpositiv, rechts positiv. Diese enge Vorzeicheninkonsistenz wird von
ATS-IV-Anwendungen und möglichen Reparaturen getrennt. Für globale lineare
Größenfamilien mit definierten Referenzsummen liefert Starrheit höchstens
einen $y$-abhängigen Skalar; klassische Höhe und klassischer Grad folgen
nur mit den angegebenen Normierungs- und Domänenidentifikationen.
Klassenmitgliedschaft und Nichtmitgliedschaft von Maßen und Log-Volumina
sind nicht etabliert. Keine Aussage über die Wahrheit von abc, die
Korrektheit von IUT oder die gesamte ATS-Kette wird behauptet.
\end{abstract}

\clearpage
\tableofcontents
\clearpage

%=====================================================================
\section{Einleitung}\label{sec:intro}

\subsection{Mandat und Ort in der Serie}

Diese Serie untersucht die publizierte Landschaft um die abc-Vermutung mit
einem festen Auditierungsinstrument. Teil 1~\cite{part1} lokalisierte den
umstrittenen Schritt von Mochizukis Korollar
3.12~\cite{mochizukiIUT,scholzestix}; die Teile 2--3 entwickelten das
Instrument: Teil 2 führte das \emph{Brückenvertrag}-Kriterium
ein~\cite{part2}, und Teil 3 dessen Zertifikatsform~\cite{part3} (den
v1-F-Standard), die im Voraus festlegen, was ein Transportargument zwischen
zwei Rahmenwerken liefern muss, bevor seine Schlussfolgerung verwendet werden
darf; Teil 4~\cite{part4} wandte das Zertifikat auf Kirti Joshis
ATS-Serie als \emph{Form-Audit} an: ein vorab festgelegter vierzeiliger
Test, ob die Kette ihr Transportgesetz, ihren Kompensationsmechanismus,
ihren zweiten Modus und ihre container-übergreifende Vergleichbarkeit
explizit genug angibt, um überhaupt auditierbar zu sein. Das Ergebnis von
Teil 4 war, bewusst, eine Aussage über Auditierbarkeit, nicht über Substanz.
Ein Werkzeug-Supplement analysiert die Normierungsidentifikation
gesondert~\cite{gnc}, und ein begleitendes Landschaftspapier verortet die
Serie unter den Wegen des weiteren Programms~\cite{landscapeA}.

Der vorliegende Teil 5 ist ein \emph{gezielter} Substanz-Audit ausgewählter
tragender interner Schritte --- keine vollständige Verifikation der Kette.
Sein vor Arbeitsbeginn festgelegtes Mandat ist die von Teil 4 offen
gelassene Frage: Tragen die ausgewählten tragenden Schritte der publizierten
Kette --- das heißt, etablieren die in den zitierten Versionen
geschriebenen Beweise die Aussagen, die die Kette an diesen Stellen
verwendet? IUT-Importe, die Existenzbasis und Schritte außerhalb der
ausgewählten Stellen liegen außerhalb des Audits.

\subsection{Methode}

Der Audit folgte fünf Phasen, die vor der Syntheseetappe in einem
zeitgleich geführten internen Beweisregister festgehalten wurden. Phase 0
kartierte den \emph{Beweis-Spine}: die minimale Kette von Resultaten, die
von der abc-artigen Ungleichung zurück zur Existenz des Deformationsraums
führt, mit einer Import-Triage, die die eigenen Resultate der Kette von den
aus dem IUT-Korpus zitierten Aussagen trennt. Die Phasen 1--3 verifizierten
dann, zeilenverankert gegen die in Abschnitt~\ref{sec:findings} genannten
exakten arXiv-Versionen, die vom Spine als stark tragend priorisierten
Stellen: den Volltext des Normierungsmechanismus, das Basispaket
(Nichtkonstanz-Satz, Renormierungs-Korollar, adelischer
Zusammensetzungssatz) und den Kern der unteren Schranke samt seiner
Korollar-Übersetzung. Phase 4 synthetisierte die Befunde zu den hier
berichteten typisierten Befunden; die von ihr offen gelassene Lesartenfrage
wird durch die Starrheitsklassifikation von Teil 6 adressiert, und
Abschnitt~\ref{sec:case} berichtet diese Klassifikation unter ihren
benannten Hypothesen.

Drei Disziplinen aus den früheren Teilen gelten unverändert weiter. Erstens
\emph{wörtliche Verankerung}: Jeder Befund zitiert die exakte Version und,
wo Extrakte auf Zeilenebene vorliegen, den Zeilenbereich der Passage, die er
betrifft. Zweitens \emph{Typisierung}: Befunde werden klassifiziert als
formale Buchführung (verifiziert oder rekonstruierbar), typisierte Lücken
(der Text zeigt nicht, was verwendet wird --- mit Benennung des fehlenden
Schritts) oder dokumentierte eigene Rekonstruktionen (als solche markiert,
niemals stillschweigend ergänzt). Drittens \emph{Claim-Disziplin}: \glqq Der
Beweis, wie er geschrieben ist, zeigt dies nicht\grqq{} wird niemals zu \glqq Dies
ist falsch\grqq{} aufgewertet.

\subsection{Claim-Grenzen}\label{sec:claims-intro}

Diese Arbeit stellt den aktuellen Stand eines fortlaufenden gezielten
Audits dar. Das gleichfeldrige Starrheitsresultat ist mathematisch und an
seine vollständigen Stellen-, Feld- und Größenklassen-Hypothesen gebunden.
Die meisten Quellenbefunde betreffen fehlende Implikationen, keine
Gegenbeispiele zu einem vollständigen Satz. Gesondert identifiziert
Abschnitt~\ref{sec:s7} eine enge Inkonsistenz einer gedruckten
Vorzeichenungleichung unter expliziten endlichen Tate-Stellen- und
Normierungsannahmen. Keine dieser Befundarten entscheidet die Wahrheit von
abc, IUT oder der gesamten ATS-Kette. Keine neue Berechnung, kein
maschinengeprüfter Beweis und keine unabhängige Gesamtkettengültigkeit
werden behauptet. Quellenrevision, Reparatur oder Gegenargument können
die zulässigen Aussagen ändern, wenn sie die fehlenden Verträge liefern.

Die veröffentlichten v0.2-PDFs und der spätere hashgebundene zweisprachige
Quellenstand sind verschiedene Baselines. Zwischenarbeiten ergänzten
Teil-2-Zitation, Tabellenetiketten und eine Antwort-Routentabelle,
sprachreine Offenlegung und neues Layout. Die vorliegenden wissenschaftlichen
Änderungen verwenden die inventarisierten Koordinaten-, Feld-,
Normierungs-, Hüllen- und Vorzeichenkorrekturen; die exakte ursprüngliche
v0.2-TeX-Datei wurde nicht wiedergefunden. Diese Provenienzgrenzen machen
die Zwischenfassung nicht zu einer unabhängig verifizierten Kopie der
veröffentlichten Quelle.

\section{Der Beweis-Spine der Serie}\label{sec:spine}

Phase 0 kartierte die publizierte Kette rückwärts, von der abc-artigen
Schlussfolgerung bis zur Aussage, dass Deformationsräume der verlangten Art
existieren, anhand von Klartext-Extrakten der vier zu Beginn von
Abschnitt~\ref{sec:findings} aufgeführten Versionen. Die Karte ist
\emph{innerhalb} dieser Extrakte vollständig und abgeschnitten, nicht
geschlossen, überall dort, wo die Kette frühere Arbeiten der
Konstruktionsserie aufruft, die für den Audit nicht verfügbar waren --- am
folgenreichsten auf der Existenzebene, die auf der ersten von
ihnen~\cite{joshiI} ruht.

\subsection{Die Kette}

Tabelle~\ref{tab:spine} listet die Knoten auf, die unmittelbar auf der
Hauptkette liegen, in der Reihenfolge, in der der Beweistext sie durchläuft.
Die Bezeichnungen $S1,\dots,S22$ sind unsere eigene Buchführung; die
Satznummern sind die der Quellen.

\begin{table}[htbp]
\centering
\footnotesize
\setlength{\tabcolsep}{4pt}
\caption{Tragende Knoten, rückwärts von der abc-artigen Schlussfolgerung.
\glqq Beweis\grqq{} zählt Extraktzeilen zwischen dem \texttt{Proof.}-Marker und dem
Ende des Arguments: ein Indikator für Prüfbarkeit, keine Messung.}
\label{tab:spine}
\begin{tabularx}{\linewidth}{@{}l >{\raggedright\arraybackslash}X >{\raggedright\arraybackslash}X r@{}}
\toprule
 & Resultat & Stützt sich auf & Beweis \\
\midrule
S1  & IV, Satz 7.2.1 --- abc und arithmetisches Szpiro über $\mathbb{Q}$ & S2 & keiner \\
S2  & IV, Satz 7.1.1 --- Vojta-Ungleichung (deklarierter Hauptsatz) & S3, S9, S10 & 121 \\
S3  & IV, Satz 6.1.1 --- First Main Bound & S4 allein & 33 \\
S4  & IV, Satz 6.10.1 --- Second Main Bound (Log-Volumen-Sandwich) & S5, S6, S7, S9 & $55{+}160$ \\
S5  & die mittlere Gleichheit von S4 (ohne Satznummer) & \glqq same global normalization\grqq{} & 2 \\
S6  & IV, Prop.\ 6.10.9 / 6.10.10 / 6.10.12 --- obere Schranke & IUT 4, Thm.\ 1.10, Schritte (v)--(vii) & $2/1/2$ \\
S7  & III, Kor.\ 9.11.1.1 --- untere Schranke, übersetzt & S8, durch notationelle Übersetzung & keiner \\
S8  & III, Satz 9.11.1 --- Volumen-Untergrenze für $\widetilde{\Theta}$ & S15 sowie Hüllen- und Volumenwerkzeuge & 38 \\
S9  & IV, Satz 5.7.1 --- Existenz initialer Theta-Daten & Standardliteratur & 196 \\
S15 & III, Satz 4.2.2.1 --- $\widetilde{\Sigma}_{L'}$ stabil; $j^2$-Skalierung & ATS II, ATS II(1/2), Fargues--Fontaine & 15 \\
S16 & III, Kor.\ 4.2.2.4 --- erzwungene Renormierung & S15 und die \glqq global constraint\grqq{} & 2 \\
S17 & III, Satz 4.6.1 --- globale Periodenabbildung auf $\widetilde{\Sigma}_{L'}$ & S18, S15, S16 & 3 \\
S18 & II(1/2), Satz 5.10.1 --- Produktformel als Periodenabbildung & Eigenschaften von $Y_L$ und Arithmetikoiden & 5 \\
S20 & II(1/2), Kor.\ 4.2.5 --- globaler Frobenius auf $Y_L$ & S21 & 1 \\
S22 & II(1/2), Satz 5.5.2 --- Existenz von Arithmetikoiden & Scholze; Kedlaya--Temkin; [Joshi 2021a] & 36 \\
\bottomrule
\end{tabularx}
\end{table}

In einem Satz: S1 ruht auf der Vojta-Ungleichung S2, die auf der First Main
Bound S3 ruht, die vollständig auf dem Log-Volumen-Sandwich S4 ruht; die
\emph{untere} Hälfte dieses Sandwichs läuft über S7 zu S8, die \emph{obere}
Hälfte wird unverändert aus IUT 4 übernommen (S6), und die \emph{mittlere
Gleichheit} S5 ist es, die die beiden Hälften erstmals zu einer einzigen
Kette verbindet --- gestützt durch den Normierungsapparat S15 bis S18, der
auf der Frobenius-Konstruktion S20/S21 und der Existenzebene S22 sitzt. Mit
den nicht gezeigten Hilfsknoten hat die Karte 34 Zeilen: 27 nummerierte
Knoten und 7 Nebenzeilen, von denen 15 unmittelbar auf der Hauptkette
liegen.

\subsection{Import-Triage}

Jeder Knoten wurde nach Provenienz markiert: publizierte Standardliteratur;
Resultate, die die Serie selbst beweist; aus dem IUT-Korpus zitierte
Aussagen; und rein definitorische Passagen. Die dritte Markierung von der
ersten getrennt zu halten ist die strukturell wichtige Entscheidung der
Kartierung, denn der IUT-Korpus ist hier nicht Hintergrundliteratur, sondern
das Artefakt, das die Kette zu ersetzen antritt. Das Ledger der
Standardliteratur hat 17 Einträge, 11 davon tragend in dem Sinne, dass ein
Beweis der Kette unmittelbar auf ihnen ruht. Das Ledger der IUT-Importe hat
13 Einträge, vier davon \emph{unausgeführte Vollimporte}, bei denen der
gesamte Beweis aus einem Verweis besteht: \glqq This is the last equation on
[Mochizuki, 2021d, Step (v) on Page 658, Proof of Theorem 1.10] and its
proof is all of step (v)\grqq{} (IV:3361--3362), mit den parallelen
Formulierungen für die Schritte (vi) und (vii) bei IV:3372 und
IV:3390--3391. Diese drei Propositionen bilden die gesamte obere Hälfte des
Sandwichs S4.

Damit ist die Grenze dieses Audits festgelegt. Gemessen an der Last rangieren
die IUT-zitierten Propositionen an erster Stelle; gemessen an der Neuheit
rangieren sie an letzter, da sie die Rechnung der Quelle wörtlich sind. Ein
zeilenverankerter Audit dort wäre ein Audit von IUT 4, nicht der
untersuchten Kette, und liegt außerhalb des Mandats von
Abschnitt~\ref{sec:claims-intro}. Sie werden verzeichnet und beiseite
gelegt, nicht verifiziert.

\subsection{Beweistext an den tragenden Knoten}

Von den 34 kartierten Zeilen nehmen 21 Beweise eine von drei kurzen Formen
an: gar kein Beweisblock (3), ein Beweis, der aus einem externen Zitat
besteht (6), und ein Beweis, der aus einem Satz der Form \glqq clear\grqq{},
\glqq immediate\grqq{} oder \glqq tautology\grqq{} besteht (12). Zwölf der 21 sitzen an
Knoten, die unmittelbar auf der Hauptkette liegen: \glqq The middle equality is
a tautology using the fact that the number fields provided by $y_0'$ and
$y_0$ are identically normalized\grqq{} (S5, IV:3272--3273); \glqq This is clear from
Theorem 4.2.2.1 and the global constraint placed by the requirement that the
product formula holds\grqq{} (S16, III:1613--1614); \glqq The proof is clear\grqq{} (S20,
II(1/2):1318).

Der Befund ist deskriptiv. Ein zweizeiliger Beweis kann völlig stichhaltig
sein, und von mehreren der aufgeführten wird in
Abschnitt~\ref{sec:findings} gezeigt, dass sie genau so tragen, wie sie
geschrieben sind: Kürze indiziert hier \emph{Prüfbarkeit}, die ein Audit
adressieren kann, nicht Korrektheit, die er in dieser Form nicht adressieren
kann. Die Zählungen sind Zählungen von Extraktzeilen, in denen Seitenzahlen
als freistehende Zeilen stehen und Displays über Zeilen umbrechen, sodass
jede Längenangabe ein Indikator ist und keine unten ein Argument trägt.

\subsection{Was zur Verifikation ausgewählt wurde}

Phase 0 priorisierte die Knoten qualitativ nach getragener Last, Neuheit und
Kürze des Beweistextes, wobei die Neuheit aus dem soeben genannten Grund der
unterscheidende Faktor ist.

Die Rangfolge setzte die mittlere Gleichheit S5 an die erste Stelle: der
einzige Knoten, der die beiden Hälften zu einer Ungleichungskette verbindet,
ohne Satznummer, in den zwei oben zitierten Zeilen abgehandelt, während
dieselbe Arbeit an drei Stellen feststellt, dass lokale Vergleiche zwischen
$y_0$ und $\varphi(y_0)$ gerade nicht verfügbar sind (IV:3208--3211,
IV:772--775, IV:816--819). Teil 4 dieser Serie behandelte S5 unter der
Transportgesetz-Zeile des Brückenzertifikats und klärte seinen Status dort
--- konditional: Die Klärung gilt auf der gebundenen und der freien Lesart
der in Abschnitt~\ref{sec:fneu1} offengelegten Gewichtsfamilie und nicht auf
der stellenselektiven Lesart, und Abschnitt~\ref{sec:case} berichtet, wie
die Lesarten unter benannten Hypothesen klassifiziert werden. S5 wird daher
hier nicht erneut auditiert.

Die fünf verbleibenden gerankten Knoten folgen, eingeleitet von der
Volltextfrage, die Teil 4 offen gelassen hatte. \textbf{Der
Normierungsmechanismus} --- was \glqq normalisiertes Arithmetikoid\grqq{} operativ
bedeutet --- kommt zuerst, weil die von ihm festgelegte Lesart von den drei
nachfolgenden Knoten vorausgesetzt wird. \textbf{S18} trägt den gesamten
Produktformel-Apparat auf einem fünfzeiligen Beweis, dessen begründender
Satz besagt, dass Normierungen \emph{nicht} übertragen werden.
\textbf{S16} behauptet in zwei Zeilen, dass die $j^2$-Skalierung anderswo
eine Renormierung \emph{erzwingt} --- die Voraussetzung dafür, dass S17 und
S18 überhaupt wohldefiniert sind. \textbf{S17} wird vom Autor selbst als
\glqq vital to establishing [Mochizuki, 2021c, Corollary 3.12]\grqq{} und als der Typ
von Proposition markiert, der in den konkurrierenden Arbeiten \glqq claimed
without proof\grqq{} sei (III:1838--1840), und dann in drei Zeilen bewiesen.
\textbf{S8} ist die einzige substanziell unabhängige untere Schranke der
Kette, deren entscheidende Reduktion als klar aus der Konstruktion erklärt
wird. \textbf{S7} ist der Knoten, den ATS IV tatsächlich zitiert: kein
Beweisblock, präsentiert als notationelle Übersetzung von S8.

%=====================================================================
\section{Zeilenverankerte Befunde}\label{sec:findings}

Alle nachstehenden Befunde sind an vier Extrakte verankert; alle vier
Versionen werden in ihren Revisionen vom 24. Februar 2025 zitiert. Die
Zeilenanker beziehen sich auf unsere internen Klartext-Extrakte und sind
Findmittel, kein öffentlicher Reproduzierbarkeitsstandard; ein
Provenienz-Manifest (arXiv-Version, Extraktionswerkzeug, Prüfsummen,
Zeilennummerierungsregel) ist als Supplement der Serie geplant, und bis es
existiert, sollten die Anker zusammen mit dem zitierten Wortlaut gelesen
werden, der wörtlich ist.
\begin{itemize}
\item \textbf{IV} = ATS IV, arXiv:2403.10430v2~\cite{joshiIV};
\item \textbf{III} = ATS III, arXiv:2401.13508v4~\cite{joshiIII};
\item \textbf{II(1/2)} = ATS II(1/2), arXiv:2305.10398v12~\cite{joshiIIhalf};
\item \textbf{II} = ATS II (lokaler Prototyp), arXiv:2303.01662v3~\cite{joshiII}.
\end{itemize}
Ein Anker der Form III:6840--6843 verweist auf Zeilen des
Klartext-Extrakts, nicht auf Seiten des publizierten Dokuments. Die Extrakte
sind in einer Hinsicht verlustbehaftet, die einschränkt, was behauptet
werden darf: Der II(1/2)-Extrakt verliert griechische Buchstaben, sodass
kein Befund darauf ruht, wo ein $\lambda$ in einer Formel jener Arbeit
steht, während die III- und IV-Extrakte Tief- und Hochstellungen
(Kontrollfall III:1558) sowie die Unterscheidung $\xi$/$\Xi$ erhalten. Die vorliegende S7-Formel wurde direkt am archivierten PDF
geprüft, dessen Betragsstriche der Extrakt verlor; rekonstruierte Formeln
werden von dieser gedruckten Aussage ausdrücklich unterschieden. Eine
Grenze kehrt in den Abschnitten~\ref{sec:s8} und~\ref{sec:s7} wieder und
wird hier einmal ausgesprochen: Mochizukis Korollar 3.12 ist nicht Teil des
auditierten Korpus, sodass die Behauptung, eine Aussage der Kette
\emph{sei} jenes Korollar (III:6880--6881, IV:3197--3198), hier nicht
entscheidbar ist und nichts im Folgenden sie entscheidet.

\subsection{Der Normierungsmechanismus}\label{sec:fneu1}

Die Normierungskoordinate aus ATS II($\tfrac12$) \S\,5.3 ist der positive
Exponent, der den Untilt-Betrag in den festen Standardbetrag überführt.
Unter der punktweisen Identität (5.3.3) liefert direkte Substitution den
klassischen Grad auf $L^\times$ und die Höhe auf $\mathbb P^n(L)$. Dies
überträgt sich nicht automatisch auf $\mathbb P^n(\overline L)$ oder
Geradenbündelhöhen. Die Kompensationssprache in \S\,8 und ATS III,
Kor.~4.2.2.4, verlangt gesondert eine identifizierte Auswahlregel sowie
einen Feld- und Stellenvertrag.

Für einen festen Zahlkörper $L$, sämtliche Stellen und reelle,
$x$-unabhängige Effektiv\-gewichte, die die volle Standard-Produktformel für
jedes $x\in L^\times$ erfüllen, liefert Teil 6 über die Stellen hinweg
$\beta_{y,v}=c_y$. Ein möglicherweise von $y$ abhängiger globaler Skalar
bleibt; eine konstante globale Skalierung mit $j^2$ ist zulässig. Ausgeschlossen
sind ungleiche Endgewichte unter diesen Hypothesen, nicht jede Änderung
roher lokaler Normierungskoordinaten. Tests auf $L^\times$ an Stellen einer
Erweiterung $L'$ etablieren nicht allein den Allstellen-Satz auf $L'^\times$.

Offen bleibt, welche Regel und Domäne die operative Größe benutzt. Eine
punktweise normierte Höhe auf $\mathbb P^n(L)$, ein freier globaler Skalar,
eine Koordinatenhyperebene und ein Maßträger sind verschiedene Objekte.
Ohne die Brücken der Abschnitte~\ref{sec:s18}--\ref{sec:s8} folgt kein
Schluss über den gesamten Deformationsapparat.

\subsection{Die Periodenabbildung: ATS II(1/2), Satz 5.10.1}\label{sec:s18}

Die Quelle behauptet die Nichtkonstanz von $y\mapsto H_y$ (ATS II($\tfrac12$),
Satz~5.10.1; Extrakt II($\tfrac12$):2161--2243). Die lokale
Variationsprämisse liegt in ATS II, Satz~10.20.1(3), vor; davon sind die
folgenden Koordinaten- und Aktionsidentifikationen zu trennen.

\paragraph{Koordinaten und Buchführung.}
Rohe Logarithmen seien $u_v=\log|x_v|$, normierte Koordinaten
$z_v=\alpha_v u_v$. Der Grad der Quelle ist $\deg=\sum_v z_v$, ihre
angegebene gewichtete Hyperebene dagegen $H=\ker(\sum_v\alpha_v z_v)$.
Somit gilt im Allgemeinen
\[
 \ker\Bigl(\sum_v z_v\Bigr)\ne
 \ker\Bigl(\sum_v\alpha_v z_v\Bigr).
\]
Für $\alpha=(2,1)$ und $z=(1,-1)$ ist die erste Summe $0$, die zweite $1$.
Dies ist ein kleines Gegenbeispiel zur uneingeschränkten Kernidentifikation.
Es widerlegt weder den vollständigen Satz noch eine konsistent
formulierte Lesart in Rohkoordinaten.
Eine Rohlog-Hyperebene $\ker(\sum_v\alpha_v u_v)$ ist ein gesonderter
konsistenter Vertrag; sie darf nicht stillschweigend mit der doppelt
gewichteten identifiziert werden. Die Trägerbuchführung der direkten Summe
bleibt in ihrer definierten Domäne erhalten; die frühere pauschale Aussage,
dass Teile (1)--(5) tragen, wird zurückgenommen.

\paragraph{Translationen und Periodenbild.}
Ein linearer Unterraum $H$ ist genau dann unter Translation um $t$
invariant, wenn $t\in H$. Hier muss $\log_y(x)$ in der tatsächlich gewählten
Hyperebene liegen. Die Produktformel für $\sum z_v$ allein etabliert nicht
die Zugehörigkeit zu $\ker(\sum\alpha_v z_v)$. Eine Aktion auf
Arithmetikoiden steigt durch eine Periodenabbildung $F$ nur ab, wenn
\[
 F(y)=F(y')\ \Longrightarrow\ F(gy)=F(gy')
\]
für die betreffenden Aktionselemente $g$ gilt. Diese Wohldefiniertheit ist
eine gesonderte Quellenpflicht, keine Folge lokaler Normwertvariation.

\paragraph{Nichtproportionalität und verbleibende Rekonstruktion.}
Bei einem konsistent definierten gewichteten Kern lässt die Multiplikation
aller Koeffizienten mit demselben von null verschiedenen Skalar die
Hyperebene unverändert. Die Nichtkonstanz von $\lambda_y$ allein trennt
daher keine Hyperebenen; benötigt wird ein Paar nichtproportionaler
Koeffizientenfamilien. Rekonstruktion R-1 bleibt an H-REC gebunden: Der
variierte nicht-archimedische Lokalpunkt lässt sich bei festgehaltenen
übrigen Koordinaten einsetzen. ATS II, Satz~10.20.1(3), liefert die lokale
Variationsprämisse, nicht diese Rekombination oder die Periodenbildaktion.
Die archimedische Reichweite bleibt ungeprüft. Kein unbedingter
S18-Abschluss wird behauptet.

\subsection{Die erzwungene Renormierung: ATS III, Korollar 4.2.2.4}\label{sec:s16}

ATS III, Kor.~4.2.2.4 (Extrakt III:1600--1614), behauptet, dass die
$j^2$-Skalierung an $V^{\mathrm{odd,ss}}$ eine geeignete Renormierung an
den übrigen Stellen erzwingt. Die Menge betrifft ungerade Restcharakteristik
mit schlechter halbstabiler, split-multiplikativer Reduktion (ATS III
\S\,3.2(6)), nicht supersinguläre Reduktion. Notwendigkeit, Lokalisierung,
Existenz und Eindeutigkeit sind getrennte Aussagen.

\paragraph{Ein baseline\-relativer Notwendigkeitstest (R-2).}
Sei $S$ eine nichtleere endliche Menge nicht-archimedischer Stellen von
$L'$, seien die getesteten globalen Elemente in $L^\times$ und die
Stellen auf die angegebenen üblichen Beträge auf $L$ eingeschränkt.
Relativ zur zusätzlichen Basislinie unveränderter Normierungskoordinaten
außerhalb $S$ hat das skalierte Produkt den Faktor
\[
 \Bigl(\prod_{w\in S}|x|_w^{\mathrm{std}}\Bigr)^{j^2-1}.
\]
Man wähle eine rationale Primzahl $r$ unter einem $w_0\in S$, also
$r\in\mathbb Q^\times\subset L^\times$. Dann ist $\log|r|_w\le0$ auf
$S$ und an $w_0$ strikt negativ. Für $j\ge2$ ist der angegebene Faktor
von $1$ verschieden. Dies ist ein globales Element der richtigen
Testdomäne, kein beliebiger lokaler Uniformisierer und kein möglicherweise
nur in $L'$ liegender Hauptidealerzeuger. Nichtleere von $S$ und Basislinie
sind zusätzliche Hypothesen; keine folgt allein aus $\ell\ge5$.

\paragraph{Gleichfeldrige Starrheit und Feldabstieg (R-3).}
Sei $b_v$ die Standardmodulfamilie auf einem Zahlkörper $K$, mit
enthaltenen kanonischen lokalen Gradfaktoren, sodass
$\sum_v\log b_v(x)=0$ für jedes $x\in K^\times$. Sind
$\rho_{y,v}=b_v^{a_{y,v}}$ und $\lambda_{y,v}$ die rohe
Normierungskoordinate, so ist der Effektivkoeffizient
$\beta_{y,v}=a_{y,v}\lambda_{y,v}$. Starrheit beschränkt $\beta$, nicht
$a$ und $\lambda$ getrennt. Unter dieser Modulkonvention folgt
$\beta=c\mathbf1$. Bei rohen verlängerten Absolutbeträgen folgt
stattdessen Proportionalität zum kanonischen Gradgewichtsvektor, keine
ungewichtete Konstanz. Erforderlich sind sämtliche Stellen von $K$
einschließlich der archimedischen, reelle $x$-unabhängige Koeffizienten und
die volle gewichtete Standard-Produktformel für jedes $x\in K^\times$.
Tests nur auf $L^\times$ über $w\in V_{L'}$ implizieren nicht die
Konstanz jedes einzelnen $L'$-Gewichts. Benötigt wird eine explizite
$L'/L$-Abstiegs-/Aggregationsidentität oder die passende volle
$L'^\times$-Formel. Der arithmetische Transfer ist offen; der gleichfeldrige
Satz wird nicht bestritten.

\paragraph{Was ein tatsächlich gemeinsames Einfrieren impliziert.}
Unter diesen Gleichfeld-Identifi\-kationen und einer nichtleeren Einfrier\-menge
liefern zwei Lösungen bei demselben $j$, identischen $a_{j,w}$ und
eingefrorenen $\lambda_{j,w}$ identische eingefrorene $\beta_{j,w}$;
Konstanz liefert dann Eindeutigkeit bei Existenz. Nichtleere ist eine
angegebene Voraussetzung, keine allgemein geprüfte Quellenbehauptung.
Erfüllt die gemeinsame Referenzfamilie zusätzlich an einer eingefrorenen Stelle
\[
 \rho_{j,w}=\rho_{1,w}^{j^2},\qquad
 a_{j,w}=j^2 a_{1,w},\qquad \lambda_{j,w}=\lambda_{1,w},
\]
so folgt algebraisch
\[
 c_j=\beta_{j,w}=j^2\beta_{1,w}=j^2c_1.
\]
Diese zulässige globale Skalierung verletzt die Starrheit nicht. Ihre
Quellenanwendung verlangt weiterhin die gemeinsamen Feld-, Norm-,
Stellen- und Freeze-Identifikationen; sie liefert kein allgemeines Verdikt
projektiver Gehaltlosigkeit.

\paragraph{Der Rohkoordinaten-Freeze-Vergleich.}
Um zwei Exponententypen nicht zu vermischen, bezeichne $A_{j,v}$ den
Koeffizientenvektor im Vergleich mit festen Rohkoordinaten, nicht den
lokalen Betrags-Exponenten $a_{j,v}$ oben. Im inventarisierten Vergleich
hat er für $\varnothing\ne S\ne V$ und $j\ge2$ auf $S$ die Einträge
$\lambda_{1,v}$ und außerhalb $S$ die Einträge $j^2\lambda_{1,v}$;
unter dem zugehörigen kanonischen Vertrag ist er
$j^2\lambda_{\mathrm{can},j,v}$. Folglich
\[
 H_{\mathrm{Freeze},j}=H_{\mathrm{can},j},
 \qquad H_{\mathrm{can},j}\ \hbox{muss nicht gleich}\ H_1\ \hbox{sein}.
\]
Dieser bedingte Vergleich zertifiziert nicht die fehlende Feldidentifikation
der Quelle. Aus ungeprüften feldübergreifenden Identifikationen roher und
effektiver Vektoren folgt kein pauschaler Schluss über projektive
Information. Die wörtliche Endlichstellen-Konvention von ATS III und ihre
mögliche archimedische Vervollständigung bleiben verschiedene offene
Quellenverträge.

\subsection{Die globale Periodenabbildung: ATS III, Satz 4.6.1}\label{sec:s17}

Der Satz hat sieben Teile (III:1746--1817) und transportiert den Apparat von
Abschnitt~\ref{sec:s18} von einem einzelnen Arithmetikoid über $L$ zu einem
$\ell^*$-Tupel über $L'$, das auf Mochizukis Ansatz
$\widetilde{\Sigma}_{L'}$ eingeschränkt ist. Sein vollständiger Beweis:
\glqq This is clear from [Joshi, 2023a, Theorem 5.10.1], and the properties of
$\widetilde{\Sigma}_{L'}$ established so far and detailed in
Theorem 4.2.2.1, Corollary 4.2.2.4\grqq{} (III:1819--1820).

Die Teile (1)--(3) und (5) zitieren den komponentenweisen Lieferanten
mit separat indizierten Zielräumen (III:1749--1751). Ihr formaler Transport
steht unter den konsistenten Koordinaten- und Aktionsverträgen aus
Abschnitt~\ref{sec:s18}; das frühere unbedingte Tragen-Verdikt wird
zurückgenommen. Die Tupelerweiterung wird in einer Quellenbemerkung
(II(1/2):2267) behauptet, dort nicht bewiesen. Teil (4) beruft sich
ausdrücklich auf Korollar 4.2.2.4 (III:1791) und erbt die Feld- und
Normierungspflichten aus Abschnitt~\ref{sec:s16}. Teil (7) bleibt außerhalb
der untersuchten Passagen.

\paragraph{Der fehlende Restriktionsschritt (T17-6).} Teil (6) behauptet,
dass $z \mapsto (H_{y_1},\dots,H_{y_{\ell^*}})$ \emph{auf}
$\widetilde{\Sigma}_{L'}$ nicht konstant ist, während der Lieferant, Satz
5.10.1(7), Nichtkonstanz auf ganz $Y_L$ behauptet (II(1/2):2228--2233). Aber
$\widetilde{\Sigma}_{L'}$ ist eine echte Teilmenge des entsprechenden
Produkts, ausgeschnitten durch die scharfen Bedingungen von Definition 4.2.2
(III:1460--1467), und Nichtkonstanz auf einer Obermenge impliziert nicht
Nichtkonstanz auf einer Teilmenge. Keine der drei im Beweis zitierten
Quellen adressiert jene Passage.

\paragraph{Rekonstruktion R-4: ein bedingter Weg.}
Die Quelle liefert Stabilität von $\widetilde\Sigma_{L'}$ unter globalem
Frobenius (Satz~4.2.2.1(2), III:1536--1539) und die lokale Normrelation
$|p|_{K_{y_{n-1}}}=|p|_{K_{y_n}}^{1/p}$ (ATS~II, Satz~10.20.1(3)).
Keine dieser Aussagen beweist allein Nichtkonstanz der Periodenabbildung.
Für den vorgeschlagenen Weg ist ein gemeinsamer Koordinatenraum zu
fixieren und der tatsächliche Hyperebenen-Koeffizientenvektor $B_y$
darin zu identifizieren. Zu beweisen ist dann, dass ein einzelner
zulässiger globaler Schritt die vorgeschlagenen lokalen Änderungen
$\lambda_v\mapsto p_v\lambda_v$ gleichzeitig realisiert, dass diese die
behauptete Transformation von $B_y$ in denselben fixierten Koordinaten
induzieren und dass $B_{\varphi y}$ und $B_y$ dort tatsächlich
nichtproportional sind. Nur für nichtverschwindende Linearformen auf
diesem fixierten Raum folgt daraus $\ker B_{\varphi y}\ne\ker B_y$.
Zusammen mit Stabilität von $\widetilde\Sigma_{L'}$ und einer
wohldefinierten Periodenabbildung würde dies Nichtkonstanz auf dieser
Menge etablieren. Dies sind H-FROB-Pflichten, keine Folgen variierender
Restcharakteristiken oder roher $\lambda$-Koordinaten allein. Vorbehalte
zu Standardpunkten und archimedischen Stellen bleiben; die nötigen
globalen, Koordinaten- und Koeffizientenidentifikationen werden hier
nicht verifiziert.

\paragraph{Eine Konventionsfrage, die über Gehaltlosigkeit entscheidet.} ATS
III schreibt die Logarithmus- und Gradabbildungen der Teile (1) bis (3) ohne
die Koeffizienten $\lambda_v$ (III:1760, III:1772, III:1781), und seine
stehende Konvention erklärt alle Arithmetikoide für normalisiert. Wörtlich
gelesen ist das Funktional die Koordinatensumme und $H_{y_j}$ für jedes $z$
dieselbe Hyperebene, was Teil (6) eher ohne Inhalt als bloß unbewiesen
zurückließe; die Aussage wird erst unter der ATS-II(1/2)-Formulierung mit
den $\lambda$-Koeffizienten (II(1/2):2212--2216), die ATS III hier nicht
reproduziert, nicht gehaltlos. Dies ist ein Textbefund, kein Artefakt der
Extraktion.

\paragraph{Verdikt.} \emph{Formale Buchführung} für die Teile (1)--(3) und (5) bleibt
an die korrigierten S18-Koordinaten- und Aktionsverträge gebunden. \emph{Typisierte Lücke} bei Teil (6):
Der Beweis, wie er geschrieben ist, zeigt keine Nichtkonstanz auf
$\widetilde{\Sigma}_{L'}$; der fehlende Schritt ist die Restriktion von
$Y_L$ auf diese Teilmenge. Eine zweite, geerbte typisierte Lücke sitzt bei
Teil (4); eine dritte betrifft die Bewertungskonvention, unter deren
wörtlicher Lesart Teil (6) keinen Inhalt hätte, wobei die auditierten
Passagen nicht festlegen, welche Lesart gilt. \emph{Dokumentierte
Rekonstruktion}: R-4 ist ein konditionaler Weg: Unter H-FROB
(Abschnitt~\ref{sec:hyplabels}) etabliert es Nichtkonstanz auf
$\widetilde{\Sigma}_{L'}$, unter Verwendung von Quellen, die der Beweis
nicht zitiert; ohne eine explizite global-lokale
Verträglichkeitsaussage bleibt der Restriktionsschritt sowohl in der Quelle
als auch in unserer Rekonstruktion offen.

\subsection{Der Kern der unteren Schranke: ATS III, Satz 9.11.1}\label{sec:s8}

ATS III, Satz~9.11.1, verwendet Volumen, sein Vorbild Satz~9.9.1 dagegen
ein Supremum. Die Boxpflicht ist am tatsächlichen Träger zu prüfen;
Punktmaß null entscheidet sie nicht. Die Behauptung, Punktzugehörigkeit
könne keine Hülle positiven Maßes liefern, wird zurückgenommen.

\paragraph{Eine bedingte algebraische Hülleninklusion.}
Sei $A$ eine tatsächlich zerlegte endliche lokale Algebra,
$\mathcal O_A^{\max}$ ihre angegebene Maximalordnung und
$\mathcal O_{\mathrm{tensor}}$ die Tensorordnung mit
$\mathcal O_{\mathrm{tensor}}\subseteq\mathcal O_A^{\max}$. Es sei
vorausgesetzt, dass die tatsächliche $\operatorname{hull}(S)$ die Menge
$S$ enthält und unter Multiplikation mit $\mathcal O_A^{\max}$ stabil ist,
etwa als entsprechende Maximalordnungs-Modulhülle. Dann gilt
\[
 x\in S\ \Longrightarrow\
 x\mathcal O_{\mathrm{tensor}}\subseteq
 x\mathcal O_A^{\max}\subseteq\operatorname{hull}(S).
\]
Die letzte Inklusion verwendet diese Modulhüllen-Hypothese. Zerlegte Algebra
und Maximalordnung allein etablieren sie nicht für eine bloß konvexe oder
$\mathbb Z_p$-Modulhülle. Eine positive volldimensionale Box verlangt
zusätzlich komponentenweise Nichtverschwinden und den tatsächlichen
lokalen Träger. Keine stillschweigende Identifikation mit der ATS-Zielhülle
oder $\widetilde\Theta$ wird vorgenommen.

\paragraph{Maßnormierung und offener Volumenvertrag.}
Für eine endliche Erweiterung $E/\mathbb Q_p$, normiertes additives
Haarmaß $\mu_E(\mathcal O_E)=1$ und den Betrag
$|a|=|N_{E/\mathbb Q_p}(a)|_p^{c/[E:\mathbb Q_p]}$ mit $c>0$ und
$|p|=p^{-c}$ liefert der Variablenwechsel
\[
 \mu_E(a\mathcal O_E)=|a|^{[E:\mathbb Q_p]/c}\quad(a\in E^\times).
\]
Unter beliebiger Betragskonvention ist dies nicht allgemein $|a|$.
Das konkrete Logarithmusbild, die ATS-Zielhülle, Tensor- und
$\Gamma$-Balance, Gradgewichte, Kollation und Indexschicht sind weiterhin
zu identifizieren. Die abstrakte algebraische Inklusion liefert diesen
Volumensatz nicht.

Die im alten Vergleich verwendete endliche Exponentenidentität
$\sum_{j=1}^{n}j^2/n^2\le n$ für $n\ge1$ bleibt eine elementare
arithmetische Tatsache. Ihre Anwendung auf eine Maß-Untergrenze verlangt
die angegebenen Haar-, Grad- und Trägerverträge; das frühere pauschale
Volumenverdikt einer sauberen Exponentenarithmetik über der Lücke wird
zurückgenommen. Herkunft und Transport der $\Xi$-Objekte durch
$\log_{BK}$ bleiben ebenfalls gesonderte Quellenpflichten
(ATS III \S\S\,9.8--9.11).

\subsection{Die übersetzte untere Schranke: ATS III, Korollar 9.11.1.1}\label{sec:s7}

\paragraph{Die gedruckte Aussage.}
Das archivierte ATS-III-v4-PDF, gedruckte S.~128, Kor.~9.11.1.1, besitzt
tatsächlich Betragsstriche. In seiner Notation lautet die angegebene Formel
\[
 -\frac{1}{\ell^*}
 \left|\operatorname{LogVol}(\widetilde\Theta^{I}_{\mathrm{Mochizuki}})\right|
 \ \ge\ -\!
 \sum_{\substack{p,\ w\in V_p^{\mathrm{odd,ss}}\\
                  V_p^{\mathrm{odd,ss}}\ne\varnothing}}
 \log\left|q_w^{1/(2\ell)}\right|_{L'_w}.
\]
Die Summe läuft über die gedruckten Primzahlblöcke und ihre angegebenen
Stellen; zusätzliche reelle Gewichte stehen in dieser Formel nicht.
Der alte Extrakt verlor die Striche. Die vorgeschlagene atomare
Alternativlesart, die Behauptung erst zu ergänzender Betragsstriche und
die frühere Entlastung durch jene Alternativlesart werden ausdrücklich
zurückgenommen.

\paragraph{Enge Vorzeicheninkonsistenz mit Hypothesen.}
Es seien $\ell^*>0$, $\ell>0$, endliches reelles Log-Volumen und eine
definierte endliche, nichtleere Tate-Stellensumme mit üblicher
Tate-Betragsnormierung $0<|q_w|_{L'_w}<1$ vorausgesetzt. Jeder Summand
$\log|q_w^{1/(2\ell)}|_{L'_w}$ ist strikt negativ. Links ist der Ausdruck
nichtpositiv, rechts strikt positiv. Unter diesen expliziten Voraussetzungen
kann die gedruckte Ungleichung also nicht gelten. Dies ist eine lokale
Vorzeicheninkonsistenz dieser angegebenen Aussage, kein Verdikt über die
gesamte ATS-Kette, IUT oder abc. Leere Summen, divergente Ausdrücke oder
andere Normkonventionen werden nicht stillschweigend erfasst.

\paragraph{Was daraus nicht folgt.}
Aus einer unteren Schranke $\mathrm{Vol}\ge R$ mit $0<R<1$ folgt weder
$\mathrm{Vol}\le1$ noch $\operatorname{LogVol}\le0$; $\mathrm{Vol}=2$
ist ein Gegenbeispiel zu dieser Schlussregel. Die gesonderte lokale
Nichtpositivitätsbehauptung der Quelle darf keinen bewiesenen Volumenvertrag
ersetzen. Die Kritik an der gedruckten ATS-III-Form, ihre Verwendung in
ATS IV und eine korrigierte Untergrenzen-Aussage verlangen getrennte
Vergleiche. Eine reparierte Vorzeichenkonvention allein schließt die
Hüllen-/Maßpflichten von Abschnitt~\ref{sec:s8} nicht. H-SIGN unten hält
die tatsächlichen Betragsstriche und Voraussetzungen des Vorzeichentests
fest, keine alternative Lesart.

\section{Katalog typisierter Lücken und Verdikt-Tabelle}\label{sec:catalogue}

\subsection{Verdikt-Tabelle}

\begin{table}[htbp]\centering\small
\caption{Verdikt nach Operation; offene Brücken sind keine Gesamtkettengültigkeit.}
\label{tab:verdict}
\begin{tabularx}{\textwidth}{@{}>{\raggedright\arraybackslash}p{0.32\textwidth} >{\raggedright\arraybackslash}X@{}}
\toprule Operation & Aktuelle Reichweite \\\midrule
S18 Koordinaten/Aktion & direkte-Summen-Buchführung nur in ihrer Domäne; gewichteter Kern und Periodenbildaktion gesondert \\
S18 Nichtproportionalität & R-1 bedingt auf konsistente Koordinaten und H-REC; archimedische Reichweite offen \\
S16 Normierung & gleichfeldrige Starrheit gültig; $L'/L$, nichtleere Einfrier\-menge und Normidentifikation offen; globale $j^2$ erlaubt \\
S8 Hülle/Volumen & lokale Inklusion bedingt auf tatsächliche Modulhüllenstabilität; ATS-Haar-/Tensor-/Zielvolumenvertrag offen \\
S7 gedrucktes Korollar & Betragsstriche vorhanden; enge Vorzeicheninkonsistenz unter H-SIGN; IV und Reparaturen gesondert \\
S17 Restriktion & Restriktionsschritt der Quelle offen; R-4 bedingt auf H-FROB \\
\bottomrule\end{tabularx}\end{table}

\subsection{Dokumentierte Rekonstruktionen}

R-1 ist der Nicht\-pro\-por\-tio\-na\-li\-täts\-weg unter konsistenten Koordinaten und
H-REC; R-2 der baseline\-relative Notwendigkeitstest mit globaler rationaler
Primzahl; R-3 die gleichfeldrige Starrheit, deren Quellenanwendung
Feldabstieg verlangt; R-4 der bedingte globale Perioden-Restriktionsweg
unter dem vollständigen fixierten Koordinaten- und Koeffizientenvertrag
H-FROB. Die lokale Hülleninklusion ist an ihren tatsächlichen
Maximal\-ordnungs-Modul\-vertrag gebunden. Keine schließt die vollständige
ATS-Volumen- oder Periodenkette.

\subsection{Eine vergleichende Beobachtung}\label{sec:symmetry}

Innerhalb des Vergleichskorpus von Teil 1~\cite{part1} und des vorliegenden
Audits enthalten beide Programme einen unausgeführten Übergang vom
Enthaltenseins-Typ an einer strukturell tragenden Stelle: Das Programm, das
die Kette zu ersetzen sucht, behauptet seinen entscheidenden Schritt mit
\glqq follows formally\grqq{} (vgl.\ Teil 1), die auditierte Kette behauptet ihr
Box-Enthaltensein mit \glqq clear from the construction\grqq{}
(Abschnitt~\ref{sec:s8}). Wir behaupten hier weder Identität der Objekte
noch Äquivalenz der fehlenden Implikationen noch einen neuen Audit von IUT;
die Beobachtung ist ein Vergleich der \emph{Behauptungsform} an den
jeweiligen kritischen Stellen, nicht mehr. Der Unterschied in der Form
bleibt im Übrigen so, wie Teil 4 ihn verzeichnet hat --- die auditierte
Kette stellt einen Satz mit ausgeschriebenem Beweis und benennbaren
Reparaturpflichten auf. Die elementare Exponentenidentität ist von ihrer
weiterhin offenen Maßanwendung getrennt.
Keines der beiden Programme wird durch diese Beobachtung widerlegt.

%=====================================================================
\section{Die Starrheitsklassifikation der Normierungsgewichte}\label{sec:case}

Dieser Abschnitt wendet die Starrheitsresultate von Teil 6~\cite{part6} auf
das in Abschnitt~\ref{sec:fneu1} isolierte Normierungszentrum an. Die
nachstehenden Aussagen sind Implikationen unter explizit benannten
Hypothesen, keine paarweise disjunkten Lesarten der Quelle; die Hypothesen
werden hier einmal definiert und konsistent verwendet.

\subsection{Hypothesenetiketten}\label{sec:hyplabels}

\begin{description}
\item[H-PLACES.] Ein fester Zahlkörper $K$ und seine vollständige
Stellenmenge einschließlich der archimedischen mit identifizierter
Standardbetragsfamilie. Eine Abkürzung ist nur zulässig, wenn die
ausgelassenen Faktoren nachweislich erhalten bleiben. Tests auf $L^\times$
an $L'$-Stellen verlangen eine gesonderte Abstiegsbrücke.
\item[H-PF.] Reelle $x$-unabhängige endgültige Effektiv\-gewichte erfüllen
die volle Standard-Produktformel für jedes $x\in K^\times$ ohne
unterdrückte Stellen-, Faser- oder Korrekturterme.
\item[H-5.3.3.] Die punktweise Identität (5.3.3) gilt an jeder angegebenen
Stelle und identifiziert die Effektivkoeffizienten mit $\beta_{y,v}=1$.
Ihre Verwendung auf $L^\times$ und $\mathbb P^n(L)$ etabliert keine
Erweiterungsfeld- oder Geradenbündeltransporte.
\item[H-EFF.] Die Größe besitzt eine globale \emph{lineare} Darstellung
\[
 Q_y=\sum_v\beta_{y,v}q_v,\qquad Q_{\mathrm{ref}}=\sum_vq_v,
\]
mit $y$-unabhängigen $q_v$; beide Summen sind bei kompatiblen Index- und
Ordnungsverträgen sowie hinreichender Konvergenz für die behauptete
Skalarfaktorisierung definiert. Endliche oder absolut konvergente reelle
Summen genügen. Beliebige nichtlineare Faktorisierungen sind nicht
erfasst. Die Identifikation von $Q_{\mathrm{ref}}$ mit $Q_{\mathrm{cl}}$
ist eine zusätzliche klassische Brücke.
\item[H-ID.] Operative Höhe oder Grad ist die konkret definierte Größe
auf $\mathbb P^n(L)$ oder $L^\times$ ohne unkontrollierte zusätzliche
$y$-Daten. Algebraische Erweiterungen und Geradenbündel verlangen eigene Brücken.
\item[H-REC.] Der benötigte variierte Lokalpunkt lässt sich bei
festgehaltenen übrigen angegebenen Koordinaten zu einem zulässigen
Arithmetikoid rekombinieren; konsistente Log-/Grad-/Hyperebenen-Koordinaten
und eine Periodenbildaktion sind gesonderte Pflichten.
\item[H-FROB.] Der globale Frobenius-Schritt aus R-4 realisiert die
nötigen lokalen Transformationen gleichzeitig und bleibt in
$\widetilde\Sigma_{L'}$. Er identifiziert die tatsächlichen
nichtverschwindenden Hyperebenenformen in einem gemeinsamen fixierten
Koordinatenraum, beweist dort ihre angegebene Transformation und
tatsächliche Nichtproportionalität und liefert den wohldefinierten
Periodenabbildungsvertrag. Variation roher Normierungskoordinaten allein
genügt nicht.
\item[H-SIGN.] Die wörtlich gedruckte Betragsformel von
Kor.~9.11.1.1 wird bei $\ell,\ell^*>0$, endlichem reellem Log-Volumen
und endlicher nichtleerer üblicher Tate-Stellensumme mit
$0<|q_w|_{L'_w}<1$ dem Vorzeichentest aus Abschnitt~\ref{sec:s7}
unterzogen. Dies ist keine alternative Konvention.
\end{description}
H-5.3.3 impliziert H-PF mit $\beta=1$ im definierten gleichfeldrigen
Standardvertrag. H-PF allein liefert nur $\beta=c_y\mathbf1$, keine
$y$-Unabhängigkeit. H-EFF betrifft eine definierte lineare Größenfamilie,
keine Feld- oder Stellenidentifikation; Maße werden nicht automatisch
klassifiziert.

\subsection{Die Lesarten der Gewichtsfamilie}

Die Normierungsanalyse von Abschnitt~\ref{sec:fneu1} identifizierte drei
Lesarten der Gewichtsfamilie $(\lambda_{y,v})$ im auditierten Text und fand,
dass der Text zwischen ihnen nicht auswählt: eine \emph{gebundene} Lesart,
auf der die Normierungskoordinate durch die
Artin--Whaples-Standardnormierungen~\cite{artinwhaples} bestimmt ist
(\glqq \emph{die} Normierungskoordinate\grqq{} aus \S\,5.3); eine \emph{freie} Lesart,
auf der die Gewichte irgendeine reelle Familie sind, die die gewichtete
Produktformel erfüllt; und eine \emph{stellenselektive} Lesart, auf der eine
Regel echt stellenselektive Gewichte zuweist, während die Produktformel
weiterhin gilt --- wobei Letzteres das ist, was die \emph{Verwendung} der
Gewichte in der Kette verlangt. Diese Lesarten sind keine disjunkten
Lösungsklassen: Das gebundene Tupel ist selbst eine Produktformel-Lösung,
und eine stellenselektive Lösung wäre, falls es sie gäbe, ebenfalls eine.
Sie sind Alternativen dafür, was die Quelle \emph{meint}; welche von ihnen
die Quelle beabsichtigt, ist eine Frage der Textinterpretation, die die
nachstehende Mathematik nicht entscheidet. Was die Mathematik entscheidet,
ist, welche dieser Lesarten unter den benannten Hypothesen durch
Exponentenfamilien \emph{realisierbar} sind.

\subsection{Anwendbarkeits-Gate: Stellenbereich}\label{sec:gate}

Der Starrheitssatz von Teil 6 gilt für eine volle Allstellen-Produktformel.
Bevor er auf die auditierten Passagen angewandt werden kann, muss der in
Abschnitt~\ref{sec:s16} verzeichnete Stellenbereichs-Konflikt ausgetragen
werden: ATS II($\tfrac12$) führt $V_L$ als alle Stellen ein, archimedische
eingeschlossen, während ATS III \S\,3.2(5), wörtlich gelesen, nur
nicht-archimedische Bewertungen indiziert und seine Produktformel über
diesen Bereich laufen lässt. Der auditierte Text gleicht diese Konventionen
nicht ab, und die Potenzklassen-Äquivalenz der lokalen Beträge ersetzt die
fehlenden Stellen nicht. Daher gilt: (i) Unterdrückt ATS III die
archimedischen Faktoren, behält sie aber bei und liefert den vollen
gleichfeldrigen Testbereich sowie die übrigen H-PF-Voraussetzungen, so
greift der Starrheitssatz auf den vervollständigten Effektivgewichts-Vektor;
(ii) ist die Produktformel von ATS III wörtlich eine rein
nicht-archimedische Bedingung, so greift der Satz nicht wie ausgesprochen,
und der Status jener Passage bleibt offen, bis ein gesonderter Satz oder
eine explizite archimedische Vervollständigung vorliegt. Im Spezialfall, in
dem die Kompensation an den endlichen Stellen mit festen archimedischen
Standardgewichten kombiniert wird, fällt der kombinierte Allstellen-Vektor
unter H-PLACES und H-PF, und die Starrheit zwingt die endlichen Gewichte
dann auf einen einzigen globalen Wert; diese kombinierte Formulierung muss
jedoch von der Quelle ausgesprochen und nicht aus der Notation erschlossen
werden. Alles Weitere in diesem Abschnitt wird unter H-PLACES behauptet.

\subsection{Die Klassifikation aus Teil 6}

Teil 6 v0.3~\cite{part6} liefert gleichfeldrige Starrheit und
unveränderte endliche Faserumverteilung. Unter H-PLACES und H-PF gilt
$\beta_{y,v}=c_y$. Unter H-EFF mit seinem tatsächlichen Konvergenzvertrag
folgt
\[
 Q_y=c_y Q_{\mathrm{ref}}.
\]
Erst mit der zusätzlichen Identifikation $Q_{\mathrm{ref}}=Q_{\mathrm{cl}}$
darf dies als $c_yQ_{\mathrm{cl}}$ geschrieben werden. Unter H-5.3.3 sind
die Koeffizienten $1$; direkte Substitution liefert den klassischen Grad
auf $L^\times$ und die Höhe auf $\mathbb P^n(L)$. Punkte über dem
algebraischen Abschluss und ATS-III-Geradenbündelhöhen erben dies nicht
allein durch Notation.

\begin{table}[htbp]\centering\small
\caption{Starrheitsfolgen mit expliziten Domänen und Größenverträgen.}
\label{tab:rigidity-consequences}
\begin{tabularx}{\textwidth}{@{}>{\raggedright\arraybackslash}p{0.32\textwidth} >{\raggedright\arraybackslash}X@{}}
\toprule Vertrag & Zulässige Schlussfolgerung \\\midrule
H-PLACES, H-PF & konstante Endgewichte; globale Skalierung mit $j^2$ zulässig \\
zusätzlich H-EFF & $Q_y=c_yQ_{\mathrm{ref}}$; klassische Identifikation gesondert \\
H-5.3.3 und H-ID & klassischer Grad auf $L^\times$, Höhe auf $\mathbb P^n(L)$ \\
Maße, Hüllen, Log-Volumen & Mitgliedschaft und Nichtmitgliedschaft in H-EFF unklassifiziert; kein Starrheitsschluss \\
\bottomrule\end{tabularx}\end{table}

Für einen stabilisierten Ausdruck auf einer angegebenen nichtleeren
$L^\times$-Familie liefert H-5.3.3 denselben klassischen Wert jedes
definierten Bestandteils nur unter der angegebenen Quellennormierung;
sein Supremum hat dann diesen Wert. Andere Domänen oder zusätzliche
Trägerdaten verlangen eigene Brücken. Ein allgemeines Supremum wird
nicht stillschweigend als lineare H-EFF-Größe behandelt; ebenso wird
kein Skalar ohne nötige Vorzeichen- und Domänenannahmen herausgezogen.
Die Koordinatenvariation von $H_y$ allein beweist keine metrische Variation.

\subsection{Größenklassen-Brücke: was folgt und was nicht}\label{sec:bridge}

Die Größenklassen-Brücke ist eine domänensensitive Pflicht, kein
generischer Transfer von Koordinatenvariation. Ein Grad auf $L^\times$
oder eine Höhe auf $\mathbb P^n(L)$ wird nur durch die tatsächliche
H-EFF-/H-ID-Substitution klassifiziert. Ein stabilisierter Ausdruck
verlangt seinen eigenen Domänen- und Normierungsvertrag. Eine
Hyperebenenfamilie allein liefert keine metrische Variation. Für
Log-Volumen, Hülle, Haarmaß und Box-Inklusion sind weder Mitgliedschaft
noch Nichtmitgliedschaft in H-EFF hier etabliert.
\begin{table}[htbp]\centering\small
\caption{Größenklassen-Pflichten von Gewichten zu operativen Trägern.}
\label{tab:quantity-class-bridge}
\begin{tabularx}{\textwidth}{@{}>{\raggedright\arraybackslash}p{0.32\textwidth} >{\raggedright\arraybackslash}X@{}}
\toprule Träger & Benötigte Brücke \\\midrule
$L^\times$-Grad / $\mathbb P^n(L)$-Höhe & definierte lineare Substitution und klassische Referenzidentifikation \\
Stabilisierter Ausdruck & angegebene nichtleere Domäne, Normierung und nötige Supremums-/Skalarbedingungen \\
$H_y$ & konsistente Koordinaten und wohldefinierte Periodenaktion; allein kein metrischer Schluss \\
Hülle, Maß, Log-Volumen & tatsächlicher Träger, Haargrad, Tensor-/\,$\Gamma$- und Volumenvertrag; Klassenstatus unklassifiziert \\
\bottomrule\end{tabularx}\end{table}
Teil 6 schließt somit ungleiche Endgewichte unter H-PLACES/H-PF aus;
ohne identifizierte Größen- und Domänenbrücken entscheidet dies nicht die
gesamte Höhen-, Maß- oder Log-Volumen-Kette. Kein Scheitern nachgelagerter
Knoten wird allein aus der Klassifikation gefolgert.

\section{Antwort-Verifikationsvorlage und Fazit}\label{sec:template}

\subsection{Der Entscheidungsbaum}

Die Befunde dieser Arbeit betreffen den publizierten Textzustand zu den
angegebenen Auditdaten. Eine künftige Revision der ATS-Serie oder eine
Erwiderung ihres Autors ist das erwartete und willkommene nächste Ereignis;
der folgende, hier im Voraus festgelegte Entscheidungsbaum verzeichnet die
\emph{alternativen} Wege, auf denen ein solches Dokument die Befunde ändern
könnte. Die Wege sind keine kumulativen Anforderungen: Jeder Zweig gilt nur,
wenn dieser Weg gewählt wird, und eine legitime Widerlegung dieses Audits
selbst ist auf jedem Zweig ein weiteres zulässiges Ergebnis.

Tabelle~\ref{tab:response-routes} stellt die Routenlogik vor der
ausformulierten Prüfliste dar.

\begin{table}[htbp]
\centering
\small
\caption{Antwort-Routen und was jede Route etablieren müsste.}
\label{tab:response-routes}
\begin{tabularx}{\textwidth}{@{}>{\raggedright\arraybackslash}p{0.25\textwidth} >{\raggedright\arraybackslash}p{0.28\textwidth} >{\raggedright\arraybackslash}X@{}}
\toprule
Route & Was sich ändern würde & Mindestnachweis \\
\midrule
Fehlidentifikation & H-ID heftet die Klassifikation nicht mehr an die
tragende Größe & operative Größe und zusätzliche $y$-Daten identifizieren \\
Reparatur in der Exponentenklasse & die Quelle bleibt in H-EFF, repariert
aber die fehlenden Allstellen-Daten & H-PLACES klären, Auswahlregel angeben
und jeden globalen Skalar $c_y$ verfolgen \\
Träger außerhalb von H-EFF & der Starrheitssatz versperrt den gewählten
Träger nicht & Maß- oder Log-Volumen-Träger definieren und die nötige
Variation beweisen \\
Untere-Schranken-Route & Normierung ist nicht mehr der entscheidende
Schritt & volles Box-Enthaltensein beweisen oder alternativen
Untere-Schranken-Satz liefern \\
Text- und Zitationsvervollständigung & dokumentarische Abschlüsse werden
rekonstruierbar & R-1/R-4-Abschlüsse mit Quellen ausschreiben \\
\bottomrule
\end{tabularx}
\end{table}

\begin{enumerate}
\item \textbf{Fehlidentifikations-Route.} Zeigen, dass die operative
  tragende Höhe nicht die Höhe von Definition~8.2.1 ist oder dass sie
  zusätzliche $y$-Daten enthält --- also H-ID für die auditierten Passagen
  widerlegen. Die Klassifikation von Abschnitt~\ref{sec:case} heftet sich
  dann nicht mehr an die tragende Größe.
\item \textbf{Reparatur innerhalb der Exponentenklasse.} Wird dieser Weg
  gewählt: den Stellenbereich von Abschnitt~\ref{sec:gate} klären
  (H-PLACES), eine explizite Auswahlregel oder einen Beweis der
  Wahlunabhängigkeit für die erzwungene Renormierung angeben (die
  Wohldefiniertheitspflicht von Abschnitt~\ref{sec:s16}), die Regel gegen
  den eindimensionalen Lösungsraum bei $L = \mathbb{Q}$ prüfen und die
  Wirkung eines möglichen globalen Skalars $c_y$ durch jeden verwendeten
  Ausdruck verfolgen.
\item \textbf{Träger außer\-halb der Effektiv\-gewichts-Klasse.} Soll die
  $y$-Ab\-hän\-gig\-keit von einer Größe außerhalb von H-EFF getragen werden ---
  etwa auf der Ebene der Maße oder Log-Volumina ---, so ist diese Größe zu
  definieren, ihre Adresse in der Kette anzugeben und zu beweisen, dass sie
  die erforderliche Variation trägt; der Starrheitssatz versperrt diesen
  Weg nicht, und nichts in dieser Arbeit schließt ihn.
\item \textbf{Untere-Schranken-Route.} Unabhängig von der Normierung:
  entweder das in Satz~9.11.1 verwendete volle Box-Enthaltensein beweisen
  (Abschnitt~\ref{sec:s8}) oder einen alternativen Satz über die untere
  Schranke mit vollständiger Herleitung liefern.
\item \textbf{Text- und Zitationsvervollständigung.} Nach einem der
  vorstehenden Schritte: die rekonstruierbaren Abschlüsse der
  Abschnitte~\ref{sec:s18} und~\ref{sec:s17} mitsamt ihren Quellen in den
  Text schreiben. Dies ist eine Dokumentationspflicht, keine mathematische
  Vorbedingung der anderen Wege.
\end{enumerate}

\subsection{Fazit}

Der aktuelle gezielte Audit erhält die gleichfeldrige Produktformel-
Starrheit und die Trennung von Nichtkonstanz und Nichtproportionalität.
Rohe und normierte Koordinaten, Periodenbildaktion, Feld-/Stellenabstieg
und tatsächliche Freeze-Identifikationen werden getrennt. Globale
Skalierungen mit $j^2$ sind zulässig; ungleiche Endgewichte werden nur unter
vollständigen gleichfeldrigen Allstellen-Hypothesen ausgeschlossen. Eine
bedingte algebraische Modulhülleninklusion ist verfügbar; ATS-Logarithmus-,
Haargrad-, Tensor- und Zielvolumenverträge bleiben offen. Das tatsächlich
gedruckte Betragskorollar besitzt die enge Vorzeicheninkonsistenz aus
Abschnitt~\ref{sec:s7}; seine Verwendung in IV und eine reparierte Aussage
sind gesondert. Lineare Größenklassen-Faktorisierung verlangt definierte
Referenzsummen und Konvergenz, klassische Höhen ihre angegebenen Domänen.
Mitgliedschaft und Nichtmitgliedschaft von Maßen bleiben unklassifiziert.

Dies ist die aktuelle Reichweite des fortlaufenden Forschungsprojekts,
keine unabhängige Verifikation der gesamten ATS-Kette und kein Beweis oder
Gegenbeweis zu abc oder IUT. Die Antwortwege bleiben für Quellenrevision,
explizite Brücke oder Gegenargument gegen diesen Audit selbst offen.
Keine neue Berechnung und kein maschinengeprüfter Beweis werden geliefert.

%=====================================================================
\section*{KI-Offenlegung}
KI-Systeme unterstützten Quellenarbeit, mathematische Diskussion,
Textentwurf, Übersetzung und Satz. Historische adversariale Reviews und
ihre begrenzte Reichweite sind in den Forschungsakten dokumentiert.
Die vorliegende Korrekturvorbereitung verwendet die inventarisierten
Befunde und eine enge Originalseitenprüfung; sie führt kein neues
numerisches Experiment und keinen maschinengeprüften Beweis ein.
Ihre Vorbereitung behauptet keine unabhängige Abnahme dieses Kandidaten.


%=====================================================================
\begin{thebibliography}{99}
\raggedright

\bibitem{artinwhaples}
E.~Artin and G.~Whaples,
\emph{Axiomatic characterization of fields by the product formula for
valuations},
Bull.\ Amer.\ Math.\ Soc.\ \textbf{51} (1945), no.~7, 469--492.
DOI: 10.1090/S0002-9904-1945-08383-9

\bibitem{joshiI}
K.~Joshi,
\emph{Construction of Arithmetic Teichmuller Spaces I},
arXiv:2106.11452v4 [math.AG], 2021 (revised 24 February 2025).

\bibitem{joshiII}
K.~Joshi,
\emph{Construction of Arithmetic Teichmuller spaces II: Proof of a local
prototype of Mochizuki's Corollary 3.12},
arXiv:2303.01662v3 [math.AG], 2023 (revised 24 February 2025).

\bibitem{joshiIIhalf}
K.~Joshi,
\emph{Construction of Arithmetic Teichmuller Spaces II$\frac{1}{2}$:
Deformations of Number Fields},
arXiv:2305.10398v12 [math.AG], 2023 (revised 24 February 2025).

\bibitem{joshiIII}
K.~Joshi,
\emph{Construction of Arithmetic Teichmuller Spaces III: A `Rosetta Stone'
and a proof of Mochizuki's Corollary 3.12},
arXiv:2401.13508v4 [math.AG], 2024 (revised 24 February 2025).

\bibitem{joshiIV}
K.~Joshi,
\emph{Construction of Arithmetic Teichmuller Spaces IV: Proof of the
abc-conjecture},
arXiv:2403.10430v2 [math.AG], 2024 (revised 24 February 2025).

\bibitem{mochizukiIUT}
S.~Mochizuki,
\emph{Inter-universal Teichm\"{u}ller Theory I: Construction of Hodge
Theaters}, Publ.\ Res.\ Inst.\ Math.\ Sci.\ \textbf{57} (2021), no.~1/2,
3--207;
\emph{II: Hodge--Arakelov-Theoretic Evaluation}, \emph{ibid.}, 209--401;
\emph{III: Canonical Splittings of the Log-Theta-Lattice}, \emph{ibid.}, 403--626;
\emph{IV: Log-Volume Computations and Set-Theoretic Foundations},
\emph{ibid.}, 627--723.

\bibitem{scholzestix}
P.~Scholze and J.~Stix,
\emph{Why abc is still a conjecture},
unpublished manuscript, version of 16 July 2018.
Available at
\url{https://www.math.uni-bonn.de/people/scholze/WhyABCisStillaConjecture.pdf}

\bibitem{part1}
L.~Geiger,
\emph{The Gap in IUTchIII, Corollary 3.12: From Attempted Repair to
Structural Diagnosis},
Zenodo, 2026.
Concept DOI: 10.5281/zenodo.19960781
(latest version v6.3, 12 August 2026, DOI: 10.5281/zenodo.21899747).

\bibitem{part2}
L.~Geiger,
\emph{The Bridge-Contract Criterion: A Form Audit for Transport Arguments,
with an Instantiation in the abc Literature},
Zenodo, 2026.
Concept DOI: 10.5281/zenodo.21960476
(version 0.2, 16 August 2026, DOI: 10.5281/zenodo.21960477).

\bibitem{part3}
L.~Geiger,
\emph{The Bridge Certificate: A Line-Level Audit Instrument for
Inter-Theoretic Transfer Arguments, with Four Applications to the abc
Literature (DRAFT)},
Zenodo, 2026.
Concept DOI: 10.5281/zenodo.21925803
(version 0.2, 16 August 2026, DOI: 10.5281/zenodo.21960958).

\bibitem{part4}
L.~Geiger,
\emph{The Bridge Certificate Applied to the Joshi ATS Series: A Ledger
Audit},
Zenodo, 2026.
Concept DOI: 10.5281/zenodo.21951155
(version 0.2, 15 August 2026, DOI: 10.5281/zenodo.21951156).

\bibitem{part6}
L.~Geiger,
\emph{What the Product Formula Allows: A Rigidity Theorem and an
Effective-Weight Test for the ATS Height Subchain},
Zenodo, 2026.
Concept DOI: 10.5281/zenodo.21952486
(version 0.3, 2 October 2026, DOI: 10.5281/zenodo.23093243).

\bibitem{gnc}
L.~Geiger,
\emph{Global Normalization in Joshi's Arithmetic Teichm\"{u}ller Theory: What
the Identification Carries --- and What Carries the Theorem},
Zenodo, 2026.
Concept DOI: 10.5281/zenodo.21924253
(version 1.1, 13 August 2026, DOI: 10.5281/zenodo.21924254).

\bibitem{landscapeA}
L.~Geiger,
\emph{The abc Landscape: Obstructions, Failed Routes, and HCT Diagnostics
(DRAFT)},
Zenodo, 2026.
Concept DOI: 10.5281/zenodo.21916900
(latest version 1.1, 13 August 2026, DOI: 10.5281/zenodo.21924338).

\end{thebibliography}

\end{document}
